\paragraph{The \citet{dou2025} protocol, exactly.}
\label{sec:faithful}
Our replication above follows the reimplementation of \citet{esquinas2026}.
The original differs in four respects: a step size of $\alpha=0.01$ rather
than $0.05$; exploration that decays with the number of visits to the
current value, $\varepsilon=e^{-\beta t(v)}$ with $\beta=5\times10^{-7}$,
rather than with calendar time; one common 31-point price grid rather than
bins over each value's range; and a market maker that re-estimates its rule
by least squares over a rolling window. We repeat the experiment with the
first three exactly, keeping our market maker, for $4\times10^8$ periods
(100 markets; exploration has then fallen to about $e^{-20}$), and add the
placebo and the control that \citet{dou2025} themselves use: removing the
lagged price from the state. Table~\ref{tab:faithful} reports the results;
shocks are in deviation units, $0.05$, $0.25$ and $1$ corresponding to about
2\%, 10\% and 40\% of a trader's average order (theirs: 0.25\% to 15\%).

\begin{table}[t]
  \centering
  \caption{The \citet{dou2025} protocol ($\alpha=0.01$, value-specific
  exploration with $\beta=5\times10^{-7}$, common 31-point price grid,
  $\xi=500$, $\sigma_u=0.1$ unless noted; $4\times10^8$ periods, 100 markets).
  $\Delta$ profit is their normalised profitability $\Delta_C$. Shock columns:
  change in per-trader intensity one period after a noise shock of the given
  size in deviation units; rival reaction: change in the rival's intensity one
  period after a forced best-response deviation; both in \% of $\beta^N$;
  95\% CIs.}
  \label{tab:faithful}
  \resizebox{\linewidth}{!}{\begin{tabular}{lcccccc}
\toprule
Learners ($\xi=500$) & $\Delta$ profit & $\Delta$ intensity & Shock 0.05 & Shock 0.25 & Shock 1 & Rival reaction \\
\midrule
Price memory, $\gamma=0.95$ (their baseline) & $0.29 \pm 0.01$ & $0.18 \pm 0.01$ & $0.04 \pm 0.06$ & $-0.04 \pm 0.16$ & $1.93 \pm 0.30$ & $0.23 \pm 0.19$ \\
Price memory, $\gamma=0$ (myopic placebo) & $0.12 \pm 0.01$ & $0.074 \pm 0.005$ & $0.04 \pm 0.07$ & $-0.08 \pm 0.14$ & $2.22 \pm 0.26$ & $-0.04 \pm 0.10$ \\
Lagged value only, $\gamma=0.95$ (their control) & $0.24 \pm 0.01$ & $0.143 \pm 0.005$ & $0.00 \pm 0.00$ & $0.00 \pm 0.00$ & $0.00 \pm 0.00$ & $0.00 \pm 0.00$ \\
Price memory, $\gamma=0.95$, $\sigma_u=100$ & $0.37 \pm 0.01$ & $0.24 \pm 0.01$ & $0.00 \pm 0.00$ & $0.00 \pm 0.00$ & $0.00 \pm 0.00$ & $0.00 \pm 0.00$ \\
\bottomrule
\end{tabular}}
\end{table}

Three results stand out. First, collusion is lower than in our first
replication and below the $\Delta_C\approx0.75$ that \citet{dou2025} use to
calibrate their theoretical benchmark to their simulations in this regime
(their Section 5.3): $\Delta_C=\FaDeltaProfit$ (intensity: $\FaDeltaInt$).
Their own robustness grid puts this gap in perspective (their online
appendix, Figure IA.9). At $\sigma_u=0.1$ their baseline cell
($\alpha=0.01$, $\beta=5\times10^{-7}$) gives $\Delta_C=0.7$, but three of
its four neighbours ($\beta=5\times10^{-8}$ or $5\times10^{-6}$ at the same
$\alpha$, and $\alpha=0.05$ at the same $\beta$) give $0.3$, as does most of
the grid. Our faithful run ($\FaDeltaProfit$) and our earlier runs at
$\alpha=0.05$ under the first protocol ($0.31$--$0.33$) match those
neighbouring cells. The level of
collusion in this regime is thus sensitive to the hyperparameters in their
simulations too, and our replication lands where most of their grid does. Second, the
forward-looking learner ignores shocks of 2\% and 10\% of its order, and
responds to a shock of 40\% ($\FaShockLarge\pm\FaShockLargeCI\%$ of $\beta^N$),
a threshold-like pattern; but the myopic placebo, which cannot punish,
responds in the same way and by as much
($\FaMyopicShockLarge\pm\FaMyopicShockLargeCI\%$), and a forced deviation
still pays. The rival's reaction to a deviation is small and borderline
($\FaRival\pm\FaRivalCI\%$). Third, removing the lagged price lowers
collusion only from $\FaDeltaProfit$ to $\FaNoPriceDeltaProfit$, whereas
\citet{dou2025} find that it falls to zero. With heavy noise trading
($\sigma_u=100$) collusion is higher ($\FaHighDeltaProfit$) and no reaction of
any kind appears, as they report for over-pruning, although the level is
below the $0.7$ of their corresponding cell.

\paragraph{Session by session.} An average over sessions could hide a
minority of sessions that do punish. \citet{dou2025} classify each session as
price-trigger collusion if, one period after a noise shock that moves the
price by 1.2\%, both speculators' orders rise significantly (their online
appendix, Section 4.5), and report that the pattern holds across sessions in
this regime. We apply the same rule to each of $\PsBaseSessions$ retrained
markets, averaging $\PsBaseReps$ paired replays per market, for shocks that
move the price by 0.12\%, 1.2\%, 5.5\% and 7.2\% (their small, medium, large
and ultra-large deviations). No forward-looking session qualifies at any shock
size ($\PsBaseTrigMed$\% at the medium shock, $\PsBaseTrigUltra$\% at the
largest). At the medium shock $\PsBaseFlatMed$\% of sessions show no
significant change in either trader's order, and the interquartile range of
the session responses is $[\PsBaseIqrLoMed, \PsBaseIqrHiMed]$\% of the
mean order. The only sessions the rule flags belong to the myopic placebo
($\PsMyopicTrigLarge$\% at the large shock), which cannot be punishing.

We thus do not reproduce three things: the level of collusion, its
dependence on the lagged price, and a shock response that a myopic learner
does not share.

Two differences in protocol remained. The first is their stopping rule:
each session learns until its greedy strategies are unchanged for $10^6$
consecutive periods, whereas on-path strategies in our fixed-horizon runs
still change ($\FaPolicyChange$ of entries over the last $10^7$ periods). With
their rule (capped at $2\times10^9$ periods) every session converges, after
$\SrConvMin$ to $\SrConvMax\times10^8$ periods, and nothing changes:
$\Delta_C=\SrDeltaProfit$ with the lagged price and $\SrNoPriceDeltaProfit$
without it, with the same shock and deviation responses.
The second is the market maker. We replace ours by theirs, re-estimating
$\mathbb E[v\mid y]$ by least squares over the last $10^4$ periods and pricing
$p=\hat g_0+\hat\lambda y$, so that every element of their published protocol
is matched. Again all sessions meet the stopping rule (median
$\FpConvMed\times10^8$ periods), and again nothing changes: $\Delta_C=\FpDeltaProfit$
with the lagged price and $\FpNoPriceDeltaProfit$ without it; no response to a
2\% shock ($\FpShockSmall\pm\FpShockSmallCI\%$ of $\beta^N$), a response to a
40\% shock ($\FpShockLarge\pm\FpShockLargeCI\%$) of the size the myopic placebo
also shows; no significant reaction of the rival to a forced deviation
($\FpRival\pm\FpRivalCI\%$); and deviating pays.

What we cannot reproduce is therefore the evidence that separates their
price-trigger regime from over-pruning, under an independent implementation
of every element of the protocol they describe. We do not claim that
price-trigger collusion cannot arise in their implementation; we claim that
the signatures that would identify it are not robust to an independent
implementation of the same design, and that where a trigger-like response
does appear, a learner that cannot punish shows it too.
